\chapter{Introdução}\label{chap:intro}

O agronegócio brasileiro constitui um dos pilares fundamentais da economia nacional. Segundo dados do do Centro de Estudos Avançados em Economia Aplicada (CEPEA-Esalq/USP) e da Confederação da Agricultura e Pecuária do Brasil (CNA), o setor responde por aproximadamente 25\% do Produto Interno Bruto (PIB) do país e por mais de 40\% das exportações brasileiras. Nesse ecossistema, o crédito rural desempenha um papel estruturante: é o principal instrumento de financiamento da produção agropecuária, viabilizando a aquisição de insumos, máquinas e equipamentos, a implantação de culturas e a adoção de inovações tecnológicas que sustentam a competitividade do Brasil nos mercados globais de commodities agrícolas.

O marco institucional do crédito rural no Brasil remonta à Lei n.º 4.829, de 1965, que institucionalizou o crédito rural como instrumento de política pública para o desenvolvimento do setor agropecuário. Ao longo das décadas seguintes, esse sistema expandiu-se significativamente, sendo operado tanto por instituições financeiras públicas — como o Banco do Brasil, BNDES e bancos cooperativos vinculados ao Sistema OCB — quanto por bancos privados, com recursos canalizados por meio e instrumentos como o Plano Safra e demais programas governamentais de incentivo à agricultura (e.g. Programa Nacional de Fortalecimento da Agricultura Familiar — PRONAF, Moderfrota, etc).

Não obstante a relevância estratégica do setor, o ingresso de novas instituições financeiras no mercado de crédito rural apresenta desafios de ordem técnica e informacional que distinguem esse segmento do crédito corporativo e pessoal convencional. A natureza diferenciada da atividade agrícola — caracterizada por ciclos produtivos longos, sazonalidade climática, volatilidade de preços de commodities, etc — cria uma assimetria de informação, dificultando o processo de avaliação de risco e precificação do crédito. Em particular, a ausência de um portfólio histórico interno no segmento rural impede a calibração de modelos de crédito tradicionais de risco baseados em dados comportamentais próprios.

Diante desse cenário, o Sistema de Operações do Crédito Rural e do Proagro (SICOR), mantido pelo Banco Central do Brasil (BACEN), emerge como um recurso analítico viável para a construção de modelos de crédito rural. O SICOR consolida, de forma granular, informações sobre todas as operações de crédito rural contratadas pelas instituições financeiras atuantes no país, abrangendo dados sobre atividade financiada, safra, finalidade da operação, modalidade, volume contratado, inadimplência e localização geográfica do mutuário. O acesso público a esses microdados regulatórios permite que novos entrantes no mercado de crédito rural construam conhecimento analítico sobre o perfil de risco setorial sem depender de uma base histórica própria.

Este projeto propõe, portanto, o desenvolvimento de um ecossistema de Business Intelligence (BI) e Analytics e de um modelo preditivo agnóstico de risco de crédito rural, estruturado sobre os dados públicos do SICOR, para subsidiar uma instituição financeira em sua estratégia de entrada e originação ao segmento do agronegócio. A abordagem analítica proposta está fundamentada na transformação dos dados disponíveis em inteligência de negócio para tomada de decissões estruturadas, contribuindo tanto para a sustentabilidade do negócio quanto para o avanço das práticas de analytics aplicado em finanças agrícolas.
